%%%%%%%%%%%%%%%%%%%%% Chapter 27 %%%%%%%%%%%%%%%%%%%%%%%%%%%%%%%%%
% Global Compliance: GDPR, AI Act, and Vietnam PDPD
%%%%%%%%%%%%%%%%%%%%%%%%%%%%%%%%%%%%%%%%%%%%%%%%%%%%%%%%%%%%%%%%%
\motto{Compliance is not a barrier to innovation; it is the framework within which sustainable innovation is possible.}

\chapter{Global Compliance: GDPR, AI Act, and Vietnam PDPD}
\label{chap:compliance}

\abstract{The Trustworthy AI Architect operates within a complex and rapidly evolving global legal landscape. This chapter provide a comparative analysis of the three primary regulatory frameworks that shape the design of modern AI systems: the European Union's GDPR and AI Act, and Vietnam's Personal Data Protection Decree (PDPD - Decree 13/2023). We analyze the requirements for data localization, risk-based conformity assessments, and the technical fulfillment of subject rights. By mapping our architectural defenses to these legal requirements, we ensure that our systems are not only mathematically robust but also legally compliant across international borders.}

\section{The Shift Toward Rigorous Regulation}
\label{sec:compliance_shift}

In the early years of the AI revolution, innovation often outpaced governance. Today, that gap is closing. Governments around the world have recognized that AI systems, while powerful, pose significant risks to individual privacy and societal safety. For the architect, "legal compliance" is now a primary boundary condition, as essential as latency or memory constraints.

\section{Comparative Regulatory Analysis}
\label{sec:regulatory_analysis}

A Trustworthy AI Architect must navigate three primary pillars of governance:

\begin{description}
  \item[GDPR (European Union):] 
        The "gold standard" for privacy. It introduced the principles of \textbf{Privacy by Design} and the \textbf{Right to be Forgotten} (which we addressed technically in Chapter 19). It requires that any processing of personal data be necessary, proportionate, and transparent.
  \item[EU AI Act (2024):]
        The first comprehensive law dedicated specifically to Artificial Intelligence. It adopts a \textbf{Risk-Based Approach}, categorizing systems as Unacceptable, High, Limited, or Minimal risk. High-risk systems (e.g., healthcare, critical infrastructure) must undergo rigorous \textbf{Conformity Assessments} and formal verification (Chapter 22).
  \item[Decree 13/2023/ND-CP (Vietnam):]
        Vietnam's landmark Personal Data Protection Decree. It imposes strict requirements on \textbf{Data Localization} and the processing of "Sensitive Personal Data" (such as health or biometric records). It mandates that organizations perform a Personal Data Protection Impact Assessment (DPIA) and restricts the cross-border transfer of data unless specific security conditions are met.
\end{description}

\section{From Law to Architecture}
\label{sec:law_to_arch}

The Trustworthy AI Architect translates these legal prose into technical configurations. For example:
\begin{itemize}
    \item To comply with \textbf{Decree 13's} localization requirements, we use \textbf{Federated Learning} (Chapter 7) and \textbf{Small Language Models} (Chapter 17) to ensure that processing happens locally on Vietnamese soil.
    \item To satisfy the \textbf{AI Act's} transparency requirements, we implement the \textbf{Explainability Stack} (Chapter 12).
    \item To fulfill the \textbf{GDPR's} security requirements, we use \textbf{Differential Privacy} (Chapter 5) and \textbf{ZK-ML} (Chapter 21) to provide mathematical proofs of protection.
\end{itemize}

% ------------------------------------------------------------------
\section{Heuristic for Regulatory Collision}
\label{sec:regulatory_collision}
% ------------------------------------------------------------------

The comparative analysis of Section~\ref{sec:regulatory_analysis} presents each regulation as a self-consistent framework. In practice, the Trustworthy AI Architect frequently encounters a more dangerous scenario: two legally binding obligations that are \textbf{mutually contradictory} when applied to the same system. This section provides a structured heuristic for resolving the most common class of regulatory collision---the conflict between a \textbf{data subject right} (e.g., GDPR's Right to be Forgotten) and a \textbf{legal retention mandate} (e.g., a financial audit obligation).

\subsection{The Banking Agent Collision: A Running Example}
\label{sec:banking_collision}

Consider an autonomous banking agent deployed in a Vietnamese fintech that serves EU-resident customers, placing it simultaneously under the GDPR and Vietnam's Circular 09/2020 on electronic transactions, which mandates 10-year retention of all financial audit records. A customer exercises their GDPR Article 17 Right to be Forgotten. The agent's legal obligations simultaneously demand:

\begin{itemize}
  \item \textbf{GDPR Art.~17:} Erase all personal data without undue delay.
  \item \textbf{Circular 09/2020 + Basel III:} Retain all transaction records for a minimum of 10 years for financial audit purposes.
\end{itemize}

A naive ``Full Deletion'' approach violates the financial retention mandate. A naive ``No Action'' approach violates GDPR. Neither is architecturally acceptable.

\subsection{The Three-Step Collision Resolution Heuristic}
\label{sec:collision_heuristic}

\begin{description}
  \item[Step 1: Classify the Conflict Type.]
    Determine whether the collision is between (a) two \emph{data subject rights}, (b) a \emph{data subject right vs.\ a legal retention obligation}, or (c) two \emph{national sovereignty mandates} (e.g., conflicting data localization requirements). Each type requires a different resolution pattern. The banking example is Type (b).

  \item[Step 2: Apply GDPR Article 18 ``Restricted Processing'' as the Intermediate State.]
    GDPR Article 18 provides a legally recognized intermediate state between full retention and full erasure: \textbf{Restricted Processing}. Under this regime, the data is retained solely for the conflicting legal obligation, locked from all other processing (model training, analytics, profiling). The data subject is notified of the restriction, its legal basis, and the expected retention end date. This satisfies the ``legitimate interest'' exemption of GDPR Art.~17(3)(b) while honoring the data subject's intent.

  \item[Step 3: Implement the Dual-Silo Architectural Pattern.]
    The architect physically separates the restricted data from the production ML pipeline using a \textbf{Dual-Silo Architecture}: a cryptographically sealed \textbf{Legal Retention Silo} (write-once, read-only for authorized auditors, with a time-lock smart contract triggering automatic deletion upon the retention period's expiry) and an \textbf{Active ML Silo} from which the data has been fully unlearned (Chapter~\ref{chap:unlearning}). The two silos are logically and technically firewalled; a record's presence in the Legal Retention Silo does not imply its presence in the Active ML Silo's training distribution.
\end{description}

\begin{figure}[ht]
\centering
\begin{tikzpicture}[
    node distance=1.4cm and 3.0cm,
    font=\small,
    silo_box/.style={draw, rounded corners, inner sep=9pt, text centered, font=\bfseries\small, minimum height=1.2cm},
    legal_box/.style={silo_box, fill=orange!10, draw=orange!80!black},
    ml_box/.style={silo_box,    fill=blue!10,   draw=blue!80!black},
    req_box/.style={silo_box,   fill=red!10,    draw=red!80!black},
    res_box/.style={silo_box,   fill=green!10,  draw=green!80!black},
    arrow/.style={->, >=stealth, thick, draw=gray!70},
    barrow/.style={->, >=stealth, thick, dashed, draw=red!60}
]

\node (request) [req_box, text width=6cm]
    {\textbf{Deletion Request (GDPR Art.~17)}\\
     Customer exercises Right to be Forgotten};

\node (step1) [res_box, below=of request, text width=6cm]
    {\textbf{Step 1: Classify Collision (Type b)}\\
     Legal retention basis: Circular 09 / Basel III\\
     $\Rightarrow$ Restricted Processing (Art.~18)};

\node (legal) [legal_box, below=1.4cm of step1, text width=5.5cm]
    {\textbf{Legal Retention Silo}\\
     Write-once; audit-read-only\\
     Time-lock: auto-delete at year 10};

\node (ml) [ml_box, right=3cm of legal, text width=5.5cm]
    {\textbf{Active ML Silo}\\
     Machine Unlearning executed (Ch.~19)\\
     Unlearning Receipt issued};

\draw [arrow]  (request.south) -- node[right, font=\scriptsize] {Trigger collision check} (step1.north);
\draw [arrow]  (step1.south)   -- node[right, font=\scriptsize] {Lock to Legal Silo} (legal.north);
\draw [arrow]  (step1.east)    -- ++(1.5,0) |- node[right, font=\scriptsize, pos=0.25] {Unlearn from ML Silo} (ml.north);
\draw [barrow] (legal.east)    -- node[above, font=\scriptsize] {Cryptographic firewall} (ml.west);

\end{tikzpicture}
\caption{The Dual-Silo Pattern for regulatory collision resolution: a deletion request triggers Restricted Processing under GDPR Art.~18, routing data to a time-locked Legal Retention Silo, while Machine Unlearning simultaneously removes its influence from the Active ML Silo. A cryptographic firewall enforces separation.}
\label{fig:dual_silo}
\end{figure}

\subsection{Jurisdiction Priority Table for ASEAN Architects}
\label{sec:jurisdiction_priority}

Table~\ref{tab:jurisdiction_heuristic} provides a practical collision resolution reference for the most frequent regulatory pairs encountered by ASEAN-operating AI architects.

\begin{table}[ht]
\centering
\caption{Regulatory Collision Heuristic: recommended architectural resolution patterns for the most common conflict pairs in the ASEAN context.}
\label{tab:jurisdiction_heuristic}
\renewcommand{\arraystretch}{1.4}
\begin{tabular}{p{3.2cm} p{3.2cm} p{2.0cm} p{3.8cm}}
\toprule
\textbf{Obligation A} & \textbf{Obligation B} & \textbf{Type} & \textbf{Resolution Pattern} \\
\midrule
GDPR Art.~17 (Erasure) &
Financial Audit Retention &
b &
Dual-Silo + Art.~18 Restricted Processing \\

GDPR Art.~20 (Portability) &
Decree 13 Data Localization &
c &
Federated Portability (FL, no raw cross-border export) \\

EU AI Act Transparency &
Trade Secret Protection &
b &
ZK-ML Proof (proves property without revealing model weights) \\

GDPR Art.~22 (No Auto-Decision) &
SLA Latency Requirement &
b &
HOTL with async human confirmation queue \\

Decree 13 Cross-Border Restriction &
GDPR Art.~20 Portability &
c &
Sovereign Data Enclave + bilateral transfer agreement \\
\bottomrule
\end{tabular}
\end{table}

\section{Conclusion}
\label{sec:conclusion_compliance}

By aligning our architecture with global and local regulations, we build systems that are resilient to both technical and legal challenges. However, building these systems requires the right set of tools. In the next chapter, we provide a hands-on guide to the \textbf{Practitioner's Toolkit}, bridging the gap between theory and code.


